% ---------------------------------------------------------------
% CHAPTER 2: RELATED WORK
% ---------------------------------------------------------------
\chapter{Related Work}

This chapter surveys the technical domain literature that substantiated both 
principal engineering accomplishments of this internship program: the SkyFox 
Cinemas comprehensive full-stack implementation and the Robust Event Consumption 
Framework.

% ---------------------------------------------------------------
\section{Cinema Ticketing and Seat Reservation Systems}

The field of online reservation platforms has undergone extensive scholarly 
examination in the context of high-volume concurrent access systems. Established 
platforms such as BookMyShow and Ticketmaster function according to a 
\textit{broker architecture} --- acting as intermediaries facilitating connections 
between end users and cinema operators, emphasizing inventory consolidation, 
financial transaction processing, and consumer-oriented search capabilities. 
Classical research on such implementations addresses central technical requirements: 
mitigation of simultaneous double-booking and oversight of seat reservations 
during payment completion windows \cite{Bernstein2009}.

These marketplace-oriented platforms, nevertheless, do not directly service the 
operational requirements of the venue administrator. Administrative responsibilities --- 
including show rescheduling, movie assignment modifications, event cancellation, 
or in-person visitor processing --- are excluded from the scope of consumer-facing 
marketplace architectures and are conventionally managed through independent, often 
manual, administrative systems.

\textbf{How SkyFox Cinemas differs:} SkyFox has been engineered explicitly from the 
\textit{venue operator perspective}. Rather than a public-facing reservation marketplace, 
it functions as an operational control platform. The cinema administrator can specify 
movie content for given auditoriums and presentation times, alter or terminate scheduled 
performances, and reassign movie programming for any auditorium --- all accessible through 
a specialized management API. Moreover, SkyFox delivers \textit{granular visitor tracking}: 
within a unified reservation, individual patrons can be registered independently. For 
example, within a seven-seat reservation, the ticketing agent can complete check-in for 
three patrons at arrival and finalize registration for the remaining four at a later time, 
preserving individual verification status throughout the data layer. This operational 
visibility is not present in marketplace designs, which regard bookings as indivisible 
entities and lack on-site attendance monitoring.

% ---------------------------------------------------------------
\section{Full-Stack Web Frameworks and REST API Design}

The combination of a Go application server with a React user interface has become a 
prevalent methodology in modern web system construction. Go's \textit{Gin} framework 
\cite{Gin2024} facilitates development of efficient application programming interfaces 
through its minimal overhead and declarative middleware capabilities. React's component-based 
structure enables specification of user experiences that naturally correspond to 
structured API payloads.

Database administration tools including \textbf{Flyway} and \textbf{Liquibase} represent 
sector-standard solutions for managed, reversible data schema versioning in commercial 
environments. Inside the Go programming community, \textbf{Goose} fulfills an analogous 
responsibility, furnishing bidirectional schema transformation definitions with ordered 
release numbering.

\textbf{How SkyFox Cinemas applies these:} SkyFox implements the Gin + React design 
pattern to manage a realistic operational scenario --- customer seat reservations and 
theater administrator lifecycle management --- instead of a simplified reference CRUD 
solution. In contrast to entry-level Go + React instructional projects that typically 
disregard migration infrastructure, SkyFox incorporates Goose from the outset, mandating 
industrial-strength database administration from initial sprints. Furthermore, the system 
incorporates a goCD continuous integration environment via Bitbucket, furnishing a 
distribution and deployment architecture comparable to professional development organizations 
and seldom encountered in university-based project work.

% ---------------------------------------------------------------
\section{Agile Software Delivery}

Agile and Scrum frameworks constitute the predominant operational structure within 
commercial software development. The Technical Leadership role --- encompassing sprint 
organization, workload decomposition, code examination, and deployment automation supervision --- 
is documented as substantially impactful on organizational efficiency and engineering output 
quality \cite{Schwaber2020}. Specifically within novice or junior-level teams, a consistent 
technical authority facilitates team member integration, prevents architectural fragmentation, 
and sustains systematic code organization across delivery cycles.

\textbf{How SkyFox Cinemas applies this:} Unlike conventional campus-based team undertakings 
where supervisory authority fluctuates or operates informally, SkyFox preserved a consistent 
Technical Lead designation across all three sprints, deliberately mirroring professional 
Scrum configurations. This preservation of technical management provided uninterrupted 
engineering strategy, standardized code assessment criteria throughout the source control 
platform, and single-point responsibility for architectural integrity across an eleven-person 
contributor group --- yielding a consolidated, professionally-constructed codebase following 
completion rather than segregated, marginally-coordinated functions.

% ---------------------------------------------------------------
\section{Event-Driven Architectures and Fault Tolerance}

Apache Kafka \cite{Kreps2011} constitutes an enterprise-scale distributed event storage 
system predicated on an sequential, unchanging, segment-partitioned information store. This 
platform has received broad adoption in banking applications due to its exceptional aggregate 
throughput and durability assurances. It furnishes \textit{at-least-once} transmission assurance 
by specification, signifying that consuming processes may get duplicated transmissions following 
infrastructure interruptions or partition rebalancing. This requirement necessitates 
\textit{duplicate-resistant} implementation --- a characteristic whereby reprocessing an 
individual transmission generates no supplementary outcome --- as a realistic manufacturing 
necessity.

Recovery strategies employed in loosely-coupled computing environments have been extensively 
documented in literature from Amazon and Google operations research divisions \cite{Google2016}. 
Non-varying time intervals between attempts introduce the \textit{synchronized retry collapse} 
condition, whereby entire consumer populations attempt reprocessing simultaneously following 
shared infrastructure disruption. Geometric progression intervals with randomization 
($\text{baseDelay} \times 2^{n-1}$ with randomized adjustment) provides the recognized remedy. 
Message staging zones (DLQs), originally conceptualized by Hohpe and Woolf \cite{Hohpe2003}, 
redirect communications that surpass reclamation attempt thresholds into a separate data stream 
for forensic examination, preventing harmful communications from interfering with healthy 
operation. Prometheus \cite{Prometheus2024} implements the measurement infrastructure enabling 
perception of such mechanisms in manufacturing settings --- displaying attempt frequencies, 
staging queue magnitude, and consumer processing lateness.

\textbf{How the Resilient Event Consumer Framework applies these:} These techniques --- 
reprocessing, duplicate-resistance, message staging, and Prometheus measurement --- constitute 
well-recognized procedures independently, though typically necessitate custom integration effort 
by implementing organizations. The proposed framework aggregates all four methodologies into 
a unitary reusable Go compilation, specifically developed for IDFC FIRST Bank's event-based 
consuming implementations. Client divisions adopt the packaged functionality without mandated 
custom implementation of any of these techniques, yielding 87\% test compliance and confirmed 
complete data retention under transient system disruptions.

% ---------------------------------------------------------------
\section{Summary}

\begin{figure}[H]
\centering
\begin{tikzpicture}[node distance=0.5cm]
  \node[rectangle, fill=idfcblue, text=white, minimum width=12.5cm,
        minimum height=0.8cm, font=\bfseries, rounded corners=4pt]
    (title) {Related Work vs.\ This Internship's Contributions};
  \node[wideblock, below=0.5cm of title, text width=12.2cm,
        minimum height=1cm] (r1)
    {\textbf{BookMyShow / Ticketmaster (broker model):} Consumer-facing
     marketplace, atomic bookings, no operational management
     $\rightarrow$ SkyFox targets the cinema owner: show scheduling,
     cancellation, rescheduling, and partial per-seat check-in};
  \node[wideblock, below=0.4cm of r1, text width=12.2cm,
        minimum height=1cm] (r2)
    {\textbf{Gin + React + Goose / Flyway:} Standard full-stack and migration
     tooling $\rightarrow$ SkyFox applies these in a multi-role operational
     platform with CI/CD, not a tutorial-grade CRUD app};
  \node[wideblock, below=0.4cm of r2, text width=12.2cm,
        minimum height=1cm] (r3)
    {\textbf{Kafka + AWS/Netflix Retry + DLQ + Prometheus:} Individually
     well-documented patterns $\rightarrow$ This framework unifies all four
     into one Go library, eliminating per-team re-implementation};
  \node[widegreenblock, below=0.4cm of r3, text width=12.2cm,
        minimum height=1cm] (r4)
    {\textbf{Net result:} Both contributions apply established patterns in
     ways that are integrated, production-disciplined, and specific to
     IDFC FIRST Bank's engineering context};
\end{tikzpicture}
\caption{Related Work vs.\ Internship Contributions}
\end{figure}